%%% FORMATO E IDIOMA %%%
    \documentclass[letterpaper,DIV=15,12pt]{scrartcl}
    \usepackage[spanish,mexico,shorthands=off,es-lcroman]{babel}
    \usepackage{scrlayer-scrpage}
        \clearpairofpagestyles
        \ihead{\footnotesize \textit{Teoría de Conjuntos III}}
        \ohead{\footnotesize \textit{2026-2}}
        \cfoot{\normalfont\thepage}
        \addtokomafont{title}{\bfseries \rmfamily}
        \setlength{\headsep}{5pt}
        \newpairofpagestyles{beginstyle}{
            \clearpairofpagestyles
            \KOMAoptions{headsepline=false}
            \cfoot{\footnotesize \pagemark}
            \cfoot{\normalfont\thepage}
        }
        \addtokomafont{section}{\raggedleft\rmfamily}
        \addtokomafont{subsection}{\raggedleft\rmfamily}
        \addtokomafont{subsubsection}{\raggedleft\rmfamily}
        \setlength{\headsep}{8pt}
        \setlength{\footskip}{20pt}
        \setlength{\textheight}{600pt}
%%% PAQUETES %%%
    \usepackage{ifthen}
    \usepackage[shortlabels]{enumitem}
        \setenumerate[1]{label=\MakeLowercase{\roman*}), ref=\roman*}
        \setenumerate[2]{label=\MakeLowercase{\alph*}), ref=\alph*}
    \usepackage{amsmath}
    \usepackage{amsthm}\usepackage[T1]{fontenc}
    \usepackage{stix2}
    \usepackage{multicol}
	\renewcommand{\qedsymbol}{$\blacksquare$}
	\addto\captionsspanish{\renewcommand{\proofname}{\normalfont\bfseries Demostración}}	
	\newenvironment{solucion}{\renewcommand{\qedsymbol}{$\boxtimes$}\begin{proof}[\normalfont\bfseries Solución]}{\end{proof}}
		%\newtheoremstyle{manual}{3pt}{3pt}{\normalfont}{}{\bfseries}{.}{.5em}{#1#3}
		%\theoremstyle{manual}
		\newtheorem{proposicion}{Proposición}
        \newtheorem{definicion}{Definición}
        \newtheorem{teorema}{Teorema}
        \newtheorem{lema}{Lema}
        \newtheorem{corolario}{Corolario}
        \newtheorem*{ejercicio*}{Ejercicio}     
%%% C O M A N D O S %%%
    \renewcommand{\emptyset}{\varnothing}
    \newcommand{\tq}{:}
%% D O C U M E N T O %%
\begin{document}
\thispagestyle{plain}

    \begin{center}
        {\fontsize{30}{60}\rmfamily \textbf{Lema del Girasol}} \\ \vspace{.2cm} Teoría de Conjuntos III, 2026-2
    \end{center}
    \begin{flushright}
        \footnotesize \hfill Profesor: Luis Jesús Turcio Cuevas.\\
        \hfill Ayudante: Hugo Víctor García Martínez.
    \end{flushright}

    Un \textit{girasol} (\textit{o delta-sistema}) es un conjunto $A$, cuyos elementos son conjuntos finitos y existe un conjunto $R$ de manera que para cualesquiera $a,b \in A$ distintos, $a \cap b = R$. En este caso, $R$ es la \textit{raíz} del girasol.

    \begin{lema}
        Sea $\kappa > \omega$ un cardinal regular y supóngase que $A$ es un conjunto de conjuntos finitos, con $|A|=\kappa$. Entonces, existe un girasol $B \subseteq A$ de tamaño $\kappa$.
    \end{lema}
    \begin{proof}
        Se afirma que: para cada natural $m$ y conjunto $B$ de tamaño $\kappa$, si cada $b \in B$ tiene tamaño $|b|=m$, entonces $B$ contiene un girasol de tamaño $\kappa$. Procédase por inducción sobre $\omega \setminus 1$.
        \begin{enumerate}[\hspace{.8cm}]
            \item \textit{Base.} Si $B$ es conjunto de tamaño $\kappa$ y cada elemento de $B$ tiene tamaño $1$, entonces $B$ mismo es un girasol (con raíz $\emptyset$).
            \item \textit{Paso inductivo.} Sea $m \in \omega \setminus 1$ y supóngase que cada subconjunto de $A$ de tamaño $\kappa$, tal que todos sus elementos tienen tamaño $m$, contiene un girasol de tamaño $\kappa$. Sea $B$ un conjunto con $|B|=\kappa$ y supóngase que cada $b \in B$ tiene tamaño $m+1$.
            
            Para cada $x \in \bigcup A$ sea $B_x=\{ b \in B \tq x \in b \}$. Existen los siguientes dos casos:
            \begin{enumerate}
                \item Existe $x \in X$ tal que $|B_x|=\kappa$. En este caso, el conjunto $C=\{b \setminus \{x\} \tq b \in B\}$ tiene tamaño $\kappa$ y todos sus elementos tienen cardinalidad $m$, por lo que de la hipótesis de inducción, existe un girasol $C' \subseteq C$ de tamaño $\kappa$. Obsérvese que, $B'=\{c \cup \{x\} \tq c \in C\}$ es un girasol de tamaño $\kappa$ contenido en $B$.
                
                \item 
            \end{enumerate}
        \end{enumerate}
        
        
        \newpage
        Ahora, como cada $a \in A$ es finito,
        \[ A = \bigcup_{n \in \omega} \{ a \in A \tq |a|=n \} \, ; \]
        y, puesto que $|A|=\kappa>\omega$, debe existir $n \in \omega$ y tal que $D=\{a \in A \tq |a|=n\} \in [A]^\kappa$.
    \end{proof}

\end{document}